\documentclass[10pt,a4paper]{amsart}
\usepackage[margin=19mm]{geometry}
\usepackage{amsmath,amssymb,mathtools}
\usepackage[scr=rsfs,cal=euler]{mathalfa}
\usepackage{xfrac}
\usepackage[unicode,hidelinks,bookmarks=false]{hyperref}
\newcommand{\ie}{i.e.\ }
\newcommand{\cf}{cf.\ }
\newcommand{\resp}{respectively\ }
\newcommand{\Subsection}{\subsection}
\providecommand{\nobreakdash}{\nobreak\hbox{-}}
\setlength{\emergencystretch}{2em}
\multlinegap0pt
\usepackage{amssymb}
\usepackage{csquotes}
\usepackage{enumerate}
\newtheorem{proposition}{Proposition}
\newcommand{\K}{\Bbbk}
\title{\textbf{The ozone group of $U_q^+(B_2)$}}
\author{James G\'omez, Helbert Venegas}
\date{Source: arXiv:2609.04919v1 (CC BY 4.0)}
\begin{document}
\maketitle

\noindent\emph{Source and licence.} \textbf{The ozone group of $U_q^+(B_2)$}. Version arXiv:2609.04919v1 (CC BY 4.0), distributed by arXiv under the Creative Commons Attribution 4.0 International licence, http://creativecommons.org/licenses/by/4.0/, as recorded in the arXiv licence metadata for this exact version and shown on the abstract page https://arxiv.org/abs/2609.04919v1. Full licence notice: \texttt{licenses/arxiv-2609.04919-CC-BY-4.0.txt}.\par\smallskip

\noindent\textit{Editorial scope.} Counted unit: the article's proposition rigidez-ozono-UqB2 (source0.tex lines 506--517). The article's complete original proof of it is delivered with it (source0.tex lines 518--573, one \texttt{proof} environment of 56 lines). No mathematics is written, repaired or re-derived: the statement and the proof are the article's own lines, in the article's own notation, and no retained line is substituted or re-flowed.  No further source-local context is needed: every cross-reference of every retained line resolves inside the two delivered spans.  Nothing was omitted from any retained span, so the package carries no omission record.  No retained line contains a citation, so the wrapper carries no bibliography and no bibliography entry.  No retained line contains a figure environment, an image-inclusion command or drawing code, so no image is delivered and the package declares no asset dependency.  Editorial formatting only, in addition: an AMS \texttt{amsart} preamble that restates the article's own declarations and macros used by the retained lines, namely the environments \texttt{proposition} on separate counters, together with the standard packages those lines need.  The retained environments print in this document as Proposition 1 in order of appearance, whereas the article prints the same items with the numbering of its own full document.  The complete native source of the article as distributed in the arXiv e-print for this exact version is bundled unchanged as \texttt{sources/209399/source0.tex}.
\begin{proposition}
\label{rigidez-ozono-UqB2}
Let \(\varphi\in\operatorname{Aut}(U_q^+(B_2))\) such that
\(\varphi(e_1^\ell)=e_1^\ell\) and \(\varphi(e_2^\ell)=e_2^\ell\). Then
there exist \(a,b\in\mu_\ell\) such that
$$
\varphi(e_1)=ae_1,\qquad
\varphi(e_2)=be_2,\qquad
\varphi(e_3)=ab\,e_3,\qquad
\varphi(z)=ab^2z.
$$
\end{proposition}

\begin{proof}
Let \(A=U_q^+(B_2)\). Consider the standard grading on \(A\), determined by

$$
\deg(e_1)=\deg(e_2)=1,\qquad
\deg(e_3)=2,\qquad
\deg(z)=3.
$$
The defining relations are homogeneous, so \(A\) is a connected
\(\mathbb N\)-graded algebra. Moreover, \(A\) admits a presentation as an
iterated Ore extension over \(\K[z,e_3]\), and hence it is a domain.
Consequently, \(\deg(x^\ell)=\ell\deg(x)\) for every \(x\neq0\). From
\(\varphi(e_1^\ell)=e_1^\ell\), we obtain
$$
\ell\deg(\varphi(e_1))
=\deg(\varphi(e_1)^\ell)
=\deg(e_1^\ell)
=\ell,
$$
and therefore \(\deg(\varphi(e_1))=1\). Since
\(A_0=\K\) and \(A_1=\K e_1\oplus\K e_2\), there exist \(a,b,c\in\K\) such that
$$
\varphi(e_1)=ae_1+be_2+c.
$$
Comparing the degree-zero components in
\(\varphi(e_1)^\ell=e_1^\ell\), we obtain \(c^\ell=0\), and hence \(c=0\). Thus,
$$
\varphi(e_1)=ae_1+be_2.
$$
Now consider the quotient \(A/\langle e_1\rangle\). Since \(e_1=0\), the relations

$$
e_2e_1=q^{-2}e_1e_2-q^{-2}e_3,
\qquad
e_2e_3=q^2e_3e_2+z
$$
imply that \(e_3=z=0\). Hence, \(A/\langle e_1\rangle\cong\K[e_2]\). Reducing the equality
$$
(ae_1+be_2)^\ell=e_1^\ell
$$
modulo \(\langle e_1\rangle\), we obtain \(b^\ell e_2^\ell=0\). Since
\(\K[e_2]\) is a domain, \(b=0\), and therefore \(\varphi(e_1)=ae_1\).
Moreover, from \(\varphi(e_1^\ell)=e_1^\ell\), we obtain \(a^\ell=1\), that is,
\(a\in\mu_\ell\).
\noindent Analogously, from \(\varphi(e_2^\ell)=e_2^\ell\) and the quotient
\(A/\langle e_2\rangle\cong\K[e_1]\), we obtain
$$
\varphi(e_2)=be_2
$$
for some \(b\in\mu_\ell\). Finally,
$$
\varphi(e_3)=ab(e_1e_2-q^2e_2e_1)=ab\,e_3
\qquad\text{and}\qquad
\varphi(z)=ab^2(e_2e_3-q^2e_3e_2)=ab^2z.
$$
\end{proof}

\end{document}
